% Created with jtex v.1.0.21
% Created with the elsarticle-myst jtex template
% elsarticle layout selected by the `model` option (default: preprint)
\documentclass[preprint,authoryear]{elsarticle}

\usepackage{amsmath}
\usepackage{amssymb}
\usepackage{graphicx}
\usepackage{booktabs}

%%%%%%%%%%%%%%%%%%%%%%%%%%%%%%%%%%%%%%%%%%%%%%%%%%
%%%%%%%%%%%%%%%%%%%%  imports  %%%%%%%%%%%%%%%%%%%
\usepackage{natbib}
%%%%%%%%%%%%%%%%%  math commands  %%%%%%%%%%%%%%%%
\newcommand{\Ex}{\mathbf{\mathbb{E}}}
\newcommand{\DiscFac}{\beta}
\newcommand{\CRRA}{\rho}
\newcommand{\PermGroFac}{G}
\newcommand{\permShk}{\psi}
\newcommand{\AbsPatFac}{\Phi}
\newcommand{\Rfree}{\text{R}}
\newcommand{\tranShk}{\xi}
\newcommand{\hNrm}{h}
\newcommand{\hNrmOpt}{\bar{\hNrm}}
\newcommand{\tranShkMin}{\underline{\tranShk}}
\newcommand{\hNrmPes}{\underline{\hNrm}}
\newcommand{\MPC}{\boldsymbol{\kappa}}
\newcommand{\MPCmin}{\underline{\MPC}}
\newcommand{\WorstProb}{\wp}
\newcommand{\MPCmax}{\bar{\MPC}}
\newcommand{\mNrm}{m}
\newcommand{\mNrmCusp}{\mNrm^*}
\newcommand{\hNrmEx}{\Delta \hNrm}
\newcommand{\mNrmCuspEx}{\Delta \mNrm^*}
\newcommand{\mNrmMin}{\underline{\mNrm}}
\newcommand{\mNrmEx}{\Delta \mNrm}
\newcommand{\cFunc}{\mathbf{c}}
\newcommand{\cFuncReal}{\hat{\cFunc}}
\newcommand{\cLvl}{\mathbf{c}}
\newcommand{\PDV}{\operatorname{PDV}}
\newcommand{\PDVCoverc}{\mathcal{C}}
\newcommand{\cNrm}{c}
\newcommand{\uFunc}{\mathbf{u}}
\newcommand{\vFunc}{\mathbf{v}}
\newcommand{\vFuncOpt}{\bar{\vFunc}}
\newcommand{\cFuncOpt}{\bar{\cFunc}}
\newcommand{\vInv}{{\scriptsize \boldsymbol{\Lambda}}}
\newcommand{\vInvOpt}{\bar{\vInv}}
\newcommand{\vFuncReal}{\hat{\vFunc}}
\newcommand{\vInvReal}{\hat{\vInv}}
\newcommand{\vInvPes}{\underline{\vInv}}
\newcommand{\logmNrmEx}{\boldsymbol{\mu}}
\newcommand{\valModRte}{\boldsymbol{\Omega}}
\newcommand{\valModRteReal}{\hat{\valModRte}}
\newcommand{\logitValModRteReal}{\hat{\boldsymbol{X}}}
\newcommand{\Risky}{\mathbf{R}}
\newcommand{\Nrml}{\mathcal{N}}
\newcommand{\std}{\sigma}
\newcommand{\risky}{\mathbf{r}}
\newcommand{\equityPrem}{\pi}
%%%%%%%%%%%%%%%%%%%%%%%%%%%%%%%%%%%%%%%%%%%%%%%%%%


\usepackage{hyperref}
\hypersetup{colorlinks=true, linkcolor=purple, urlcolor=blue, citecolor=cyan}

% Economics Letters uses author-year citations via elsarticle-harv
\bibliographystyle{elsarticle-harv}

\journal{Economics Letters}

\begin{document}

\begin{frontmatter}

\title{Appendix to "The Method of Moderation"}

\author[1,2]{Christopher D. Carroll}\ead{ccarroll@jhu.edu}\author[1,2]{Alan Lujan\corref{cor1}}\ead{alujan@jhu.edu}\author[3]{Karsten Chipeniuk}\author[4]{Kiichi Tokuoka}\author[5]{Weifeng Wu}
\cortext[cor1]{Corresponding author}

\affiliation[1]{organization={Johns Hopkins University}, department={Krieger School of Arts and Sciences}, addressline={Baltimore, MD}}
\affiliation[2]{organization={Econ-ARK}}
\affiliation[3]{organization={Reserve Bank of New Zealand}}
\affiliation[4]{organization={Japanese Ministry of Finance}}
\affiliation[5]{organization={Fannie Mae}, addressline={Washington, DC}}

\begin{abstract}
This appendix provides detailed mathematical derivations and technical results supporting the Method of Moderation. Topics include: value function transformations and their relationship to the inverse value function; explicit formulas for minimal and maximal marginal propensities to consume; cusp point calculations for tighter upper bounds; Hermite interpolation slope formulas and MPC derivations; patience conditions ensuring well-defined solutions; and extensions to stochastic returns with explicit formulas for portfolio problems.
\end{abstract}

\begin{keyword}
Dynamic Stochastic Optimization \sep Consumption-Saving Models \sep Numerical Methods \\ \textit{JEL classification:} C63; D81; E21\end{keyword}

\end{frontmatter}

\section{Appendix: Mathematical Details}

\subsection{Patience Conditions Details}

Each patience condition from the main text controls a distinct way the problem could misbehave. The FVAC $0<\DiscFac\PermGroFac^{1 -\CRRA}\Ex[\permShk^{1 -\CRRA}]<1$ guarantees that even autarky, saving nothing and consuming income as it arrives, delivers finite expected discounted utility, so the consumer has a reason to value resources at all. The AIC $\AbsPatFac<1$ rules out indefinite deferral of consumption: under certainty the marginal utility of consuming now exceeds the discounted marginal utility of consuming later. Two further conditions bound wealth from opposite directions. The RIC $\AbsPatFac/\Rfree<1$ holds asset growth below the patience-adjusted discount rate, so wealth cannot explode; the GIC $\AbsPatFac/\PermGroFac<1$ holds consumption growth below permanent-income growth, which is what pins down a finite target wealth ratio. Finally, the FHWC $\PermGroFac/\Rfree<1$ keeps the present value of future labor income finite. Where these conditions fail, behavior changes qualitatively: when all hold, the consumer runs a buffer stock around a target wealth ratio; when the AIC fails, she borrows without limit; when the GIC fails but the RIC still holds, wealth grows without bound \citep{Carroll1997, SolvingMicroDSOPs, CarrollShanker2024}.

\subsection{Human Wealth Formulas}

The optimist's human wealth (assuming $\tranShk_{t+n}=1~\forall~n>0$) can be computed three ways: backward recursion $\hNrmOpt_{T} = 0$, $\hNrmOpt_{t} = (\PermGroFac/\Rfree)(1 + \hNrmOpt_{t+1})$; forward sum $\hNrmOpt_{t} = \sum_{n=1}^{T -t}(\PermGroFac/\Rfree)^{n}$; or infinite-horizon $\hNrmOpt = \PermGroFac/(\Rfree -\PermGroFac)$ when $\Rfree>\PermGroFac$. With $\PermGroFac=1$, $\hNrmOpt = 1/(\Rfree -1)$.

The pessimist's human wealth (assuming $\tranShk_{t+n}=\tranShkMin~\forall~n>0$) follows similarly: backward recursion $\hNrmPes_{T}=0$, $\hNrmPes_{t}=(\PermGroFac/\Rfree)(\tranShkMin + \hNrmPes_{t+1})$; forward sum $\hNrmPes_{t}=\tranShkMin\sum_{n=1}^{T -t}(\PermGroFac/\Rfree)^{n}$; or infinite-horizon $\hNrmPes=\tranShkMin\PermGroFac/(\Rfree -\PermGroFac)$. When $\tranShkMin=0$ (unemployment), $\hNrmPes=0$.

\subsection{Marginal Propensity to Consume Formulas}

The minimal MPC (perfect foresight consumer with horizon $T -t$) has three forms \citep{Carroll2001MPCBound}: backward recursion $\MPCmin_{t}=\MPCmin_{t+1}/(\MPCmin_{t+1}+\AbsPatFac/\Rfree)$ with $\MPCmin_T=1$; forward sum $\MPCmin_{t}=(\sum_{n=0}^{T -t}(\AbsPatFac/\Rfree)^{n})^{ -1}$; or infinite-horizon $\MPCmin=1 -\AbsPatFac/\Rfree = 1 -(\Rfree \DiscFac)^{1/\CRRA}/\Rfree$.

The maximal MPC \citep{CarrollToche2009} satisfies backward recursion $\MPCmax_{t}=\MPCmax_{t+1}/(\MPCmax_{t+1}+\WorstProb^{1/\CRRA}\AbsPatFac/\Rfree)$ with $\MPCmax_T=1$; forward sum $\MPCmax_{t}=(\sum_{n=0}^{T -t}(\WorstProb^{1/\CRRA}\AbsPatFac/\Rfree)^{n})^{ -1}$; or infinite-horizon $\MPCmax = 1 - \WorstProb^{1/\CRRA} (\AbsPatFac/\Rfree)$.

\subsection{Cusp Point Calculation}

The two upper bounds intersect at the cusp point $\mNrmCusp$ where

\begin{equation}
\label{eq:mNrmCuspFull}
\begin{array}{rclcll}
\bigl(\mNrmCuspEx + \hNrmEx\bigr)\,\MPCmin &= & \MPCmax\,\mNrmCuspEx & & \\
\mNrmCuspEx &= & \dfrac{\MPCmin\,\hNrmEx}{\MPCmax-\MPCmin} & & \\
\mNrmCusp &= & -\hNrmPes + \dfrac{\MPCmin\,\bigl(\hNrmOpt-\hNrmPes\bigr)}{\MPCmax-\MPCmin},
\end{array}
\end{equation}

where $\mNrmCuspEx\equiv\mNrmCusp -\mNrmMin > 0$ since $\MPCmax > \MPCmin$. For $\mNrm \in (\mNrmMin, \mNrmCusp]$, the tighter upper bound yields

\begin{equation}
\begin{array}{rcl}
\mNrmEx \MPCmin < & \cFuncReal(\mNrmMin+\mNrmEx) & < \MPCmax \mNrmEx \\
0 < & \cFuncReal(\mNrmMin+\mNrmEx) - \mNrmEx \MPCmin & < \mNrmEx(\MPCmax- \MPCmin) \\
0 < & \left(\frac{\cFuncReal(\mNrmMin+\mNrmEx) - \mNrmEx \MPCmin}{\mNrmEx(\MPCmax- \MPCmin)}\right) & < 1.
\end{array}
\end{equation}

This motivates the definition of the low-resource moderation ratio as in (\ref{eq:modRteLoTightUpBd}).

\subsection{Value Function Derivation}

Under perfect foresight, consumption grows at constant rate equal to the absolute patience factor $\AbsPatFac$: $\cLvl_{t+n}=\cLvl_{t}\AbsPatFac^{n}$. The present discounted value of consumption satisfies $\PDV_{t}^{T}(\cLvl)=\sum_{n=0}^{T -t}\DiscFac^{n}\cLvl_{t}\AbsPatFac^{n}=\cLvl_{t}\sum_{n=0}^{T -t}(\AbsPatFac/\Rfree)^{n}$, where we use $\DiscFac\AbsPatFac^{1 -\CRRA}=\AbsPatFac/\Rfree$. Dividing by consumption yields the PDV-to-consumption ratio $\PDVCoverc_{t}^{T}=\PDV_{t}^{T}(\cLvl)/\cLvl_{t}=\sum_{n=0}^{T -t}(\AbsPatFac/\Rfree)^{n}=\MPCmin_{t}^{ -1}$, which is unchanged for normalized variables. Defining $\PDVCoverc \equiv \lim_{T\to\infty} \PDVCoverc_{t}^{T}$, this yields the key identity $\PDVCoverc = \MPCmin^{ -1}$, connecting the infinite-horizon PDV-to-consumption ratio to the minimal MPC.

The optimist's value function satisfies

\begin{equation}
\begin{aligned}
\vFuncOpt_{T-1}(\mNrm_{T-1}) &\equiv  \uFunc(\cNrm_{T-1})+\DiscFac \uFunc(\cNrm_{T}) \\
&= \uFunc(\cNrm_{T-1})\left(1+\DiscFac \AbsPatFac^{1-\CRRA}\right) \\
&= \uFunc(\cNrm_{T-1})\left(1+\AbsPatFac/\Rfree\right) \\
&= \uFunc(\cNrm_{T-1})\PDVCoverc_{T-1}^{T}
\end{aligned}
\end{equation}

The infinite horizon expression becomes

\begin{equation}
\label{eq:vFuncPF}
\begin{aligned}
\vFuncOpt(\mNrm) &= \uFunc(\cFuncOpt(\mNrm))\PDVCoverc \\
&= \uFunc(\cFuncOpt(\mNrm))\MPCmin^{-1} \\
&= \uFunc((\mNrmEx+\hNrmEx)\MPCmin) \MPCmin^{-1} \\
&= \uFunc(\mNrmEx+\hNrmEx)\MPCmin^{-\CRRA}.
\end{aligned}
\end{equation}

This can be transformed as

\begin{equation}
\begin{aligned}
\vInvOpt &\equiv  \left((1-\CRRA)\vFuncOpt\right)^{1/(1-\CRRA)}   \\
&= \cNrm\,\PDVCoverc^{1/(1-\CRRA)} \\
&= (\mNrmEx+\hNrmEx)\MPCmin^{-\CRRA/(1-\CRRA)}.
\end{aligned}
\end{equation}

For the realist's problem, we define $\vInvReal = \left((1 -\CRRA)\vFuncReal(\mNrm)\right)^{1/(1 -\CRRA)}$. Using the bounds $\vInvPes < \vInvReal < \vInvOpt$, we define

\begin{equation}
\label{eq:valModRteReal}
\valModRteReal(\logmNrmEx) = \left(\frac{\vInvReal(\mNrmMin+e^{\logmNrmEx})-\vInvPes(\mNrmMin+e^{\logmNrmEx})}{\hNrmEx \MPCmin \,\PDVCoverc^{1/(1-\CRRA)}}\right)
\end{equation}

and the logit-transformed counterpart:

\begin{equation}
\label{eq:ChiUpper}
\begin{aligned}
\logitValModRteReal(\logmNrmEx) &= \log \left(\frac{\valModRteReal(\logmNrmEx)}{1-\valModRteReal(\logmNrmEx)}\right) \\
&= \log(\valModRteReal(\logmNrmEx)) - \log(1-\valModRteReal(\logmNrmEx))
\end{aligned}
\end{equation}

Inverting these approximations yields

\begin{equation}
\label{eq:vInvHi}
\vInvReal = \vInvPes+\overbrace{\left(\frac{1}{1+\exp(-\logitValModRteReal)}\right)}^{=\valModRteReal} \hNrmEx \MPCmin \,\PDVCoverc^{1/(1-\CRRA) }
\end{equation}

from which the value function approximation is $\vFuncReal = \uFunc(\vInvReal)$.

\subsection{I.I.D. Stochastic Returns: MPC derivation}\label{stochastic-returns-mgf-derivation}

The fact that a linear consumption function with an MPC $= 1 - (\DiscFac \Ex[\Risky^{1 -\CRRA}])^{1/\CRRA}$ satisfies the Euler equation with i.i.d. returns and no labor income can be derived by the method of undetermined coefficients.  In particular, assume that $\cFuncOpt(\mNrm) = \mNrm\MPCmin$, with a time-independent MPC $\MPCmin$ to be determined.  Substituting this into the Euler equation, we have

\begin{equation}
\label{eq:stochReturnsEulerEqn}
\begin{aligned}
1 &= \DiscFac \Ex_t\left[\Risky_{t+1} \left(\frac{\cNrm_{t+1}}{\cNrm_t}\right)^{-\CRRA}\right]\\
&= \DiscFac \Ex_t\left[\Risky_{t+1} \left(\frac{\mNrm_{t+1}}{\mNrm_t}\right)^{-\CRRA}\right]
\end{aligned}
\end{equation}

where the second equality uses the assumed form of the consumption function.  Since there is no labor income, $\mNrm_{t+1} = \Risky_{t+1}(\mNrm_t - \cNrm_t)$.  Substituting this into the above we obtain

\begin{equation}
\label{eq:stochReturnsEulerEqnContd}
1 = \DiscFac \Ex_t\left[\Risky_{t+1} \left(\Risky_{t+1}(1-\MPCmin)\right)^{-\CRRA}\right]
\end{equation}
Solving for $\MPCmin$ and recalling that returns are i.i.d. gives $\MPCmin=1 - (\DiscFac \Ex[\Risky^{1 -\CRRA}])^{1/\CRRA}$.

In the particular case of lognormal returns, the MPC can be written in closed form.  The moment generating function (MGF) for lognormal returns provides the key formula. For $\log \Risky \sim \Nrml(\mu, \sigma^2)$, the MGF is $\Ex[e^{sX}] = \exp(\mu s + \sigma^2 s^2/2)$ where $X = \log \Risky$. Setting $s = 1 -\CRRA$ and $\mu = r + \equityPrem - \std_{\risky}^2/2$ yields\footnote{Here we can interpret $\equityPrem$ as the risk premium, that is, the additional average return from holding a risky asset compared to the risk-free rate $r$.  Adjusting the average log return by the asset volatility ensures that increasing $\std_{\risky}^2$ constitutes a mean-preserving spread of the level of return.}

\begin{equation}
\Ex[\Risky^{1-\CRRA}] = \exp\left((1-\CRRA)\left(r+\equityPrem - \frac{\std_{\risky}^2}{2}\right) + \frac{(1-\CRRA)^2\std_{\risky}^2}{2}\right).
\end{equation}

Simplifying the variance terms: $(1 -\CRRA)^2\std_{\risky}^2/2 - (1 -\CRRA)\std_{\risky}^2/2 = (1 -\CRRA)[(1 -\CRRA) -1]\std_{\risky}^2/2 = -\CRRA(1 -\CRRA)\std_{\risky}^2/2$, giving the final form

\begin{equation}
\Ex[\Risky^{1-\CRRA}] = \exp\left((1-\CRRA)\left(r+\equityPrem - \CRRA\std_{\risky}^2/2\right)\right).
\end{equation}


\bibliography{main.bib}

\end{document}
